\section[A modular hyperbola count]{A modular hyperbola count with residue restrictions}
\label{sec:hyperbola}

The large-modulus case of Proposition~\ref{prop:structured-variance}
requires a rectangle count with a congruence restriction on the unit
coordinate of a modular hyperbola.  We give the completion argument with
an explicit error, retaining multiplicities when an interval is longer
than the modulus.

\begin{lemma}[Completion with a unit residue restriction]
\label{lem:hyperbola}
Let $H,T_0,T_1$ be positive integers such that $H\ge2$, $T_0$ is
squarefree, every prime divisor of $T_0$ divides $H$, and
$(T_1,HT_0)=1$.  Put $N=HT_0$.  Let $C_0\in\Z$ satisfy $(C_0,H)=1$,
and prescribe the residue pattern
\begin{equation}
\label{eq:hyperbola-pattern}
 p_0\in(\Z/T_0\Z)^\times,\qquad k_0\in\Z/T_0\Z,\qquad
 p_1\in(\Z/T_1\Z)^\times,\qquad m_1\in\Z/T_1\Z.
\end{equation}
The residue classes modulo $1$ have their usual unique interpretation.
Let $I,J$ be bounded real intervals of lengths $X,Y_2\ge0$, with either
endpoint independently included or excluded.  For a positive integer
$l$ with $(l,HT_1)=1$, let $C_l$ be the number of integer pairs $(p,m)$
satisfying
\begin{equation}
\label{eq:hyperbola-counted-set}
\begin{gathered}
 p\in I,\qquad m\in J,\qquad l\mid p,\\
 p\equiv p_0\pmod{T_0},\qquad
 p\equiv p_1\pmod{T_1},\qquad
 m\equiv m_1\pmod{T_1},\\
 pm\equiv C_0+Hk_0\pmod N.
\end{gathered}
\end{equation}
Here $p$ is an integer variable, without a primality condition.
Uniformly in all these data,
\begin{equation}
\label{eq:hyperbola-explicit-count}
 C_l=
 \frac{XY_2\varphi(N)}{lT_1^2N^2\varphi(T_0)}
 +O\left(
 \tau(N)^2\log^2(2N)\sqrt{NT_0}
 \left(1+\frac{X}{lN}+\frac{Y_2}{N}\right)
 \right),
\end{equation}
with an absolute implied constant.

In particular, suppose that $R\to\infty$ and the data range uniformly
over families satisfying
\begin{equation}
\label{eq:hyperbola-application-range}
 R^{4/5}<H\le R\mathcal Z^{11},\qquad
 X\ll R,\qquad Y_2\le3H,\qquad
 \mathcal Z,T_0,T_1=R^{o(1)}.
\end{equation}
Then, uniformly over all the patterns and intervals above and all
squarefree $l\le R^{1/10}$ with $(l,HT_1)=1$,
\begin{equation}
\label{eq:hyperbola-application-count}
 C_l=
 \frac{XY_2\varphi(N)}{lT_1^2N^2\varphi(T_0)}
 +O\bigl(R^{3/5+o(1)}\bigr).
\end{equation}
Thus the same assertion holds with error $O(R^{7/10+o(1)})$.
Uniformity in the last two assertions means that the upper bounds
represented by $R^{o(1)}$ in
\eqref{eq:hyperbola-application-range} hold uniformly over the family;
equivalently, for every $\delta>0$ the error in
\eqref{eq:hyperbola-application-count} is $O_\delta(R^{3/5+\delta})$
for sufficiently large $R$ in that family.  No lower bound for $X/l$
is required.
\end{lemma}

\begin{proof}
Write $e_n(t)=\exp(2\pi i t/n)$.  We use the classical complete
Kloosterman estimate, in the form recorded in
\cite[Lemma~7.1]{Lichtman2022}:
\begin{equation}
\label{eq:hyperbola-weil}
 \Kl(a,b;n):=\sum_{\substack{u\bmod n\\(u,n)=1}}
                  e_n(au+b\bar u),
 \qquad
 |\Kl(a,b;n)|\ll\tau(n)n^{1/2}(a,b,n)^{1/2}.
\end{equation}
This bound permits arbitrary composite $n$ and arbitrary integer
frequencies, including zero frequencies.  The inverse $\bar u$ is
taken modulo $n$.

\paragraph{Changing the first coordinate.}
Choose a representative of $k_0$ and put $c=C_0+Hk_0$ modulo $N$.
The integers $H$ and $N$ have the same prime support.  Since
$(C_0,H)=1$, it follows that $(c,N)=1$.  The hypothesis on $l$ also
implies $(l,NT_1)=1$.  On setting $p=lv$, define
\[
 \begin{gathered}
 U=X/l,\qquad V=Y_2,\qquad c'=c\bar l\pmod N,\\
 a_0=p_0\bar l\pmod{T_0},\qquad
 a_1=p_1\bar l\pmod{T_1}.
 \end{gathered}
\]
Every inverse here is taken in its indicated modulus.  The interval
$I_l=\{x/l:x\in I\}$ has length $U$ and the induced endpoint
conventions.  The pairs now satisfy
\[
 \begin{gathered}
 v\in I_l,\quad v\equiv a_1\pmod{T_1},\qquad
 m\in J,\quad m\equiv m_1\pmod{T_1},\\
 v\equiv a_0\pmod{T_0},\qquad vm\equiv c'\pmod N.
 \end{gathered}
\]
In particular $v$ and $m$ are automatically units modulo $N$.

\paragraph{Fourier weights and their divisor-weighted sums.}
For $h,k\bmod N$, define the unnormalized weights
\begin{equation}
\label{eq:hyperbola-fourier-weights}
 A(h)=\sum_{\substack{v\in I_l\cap\Z\\v\equiv a_1\ (T_1)}}e_N(-hv),
 \qquad
 B(k)=\sum_{\substack{m\in J\cap\Z\\m\equiv m_1\ (T_1)}}e_N(-km).
\end{equation}
These are sums over the integers in the progressions, so a residue
modulo $N$ is counted with its full multiplicity.  They therefore apply
unchanged when $U$ or $V$ exceeds $N$.  Counting the progressions gives
\begin{equation}
\label{eq:hyperbola-zero-weights}
 A(0)=U/T_1+O(1),\qquad B(0)=V/T_1+O(1).
\end{equation}
For a nonzero residue $h\bmod N$, let
$|h|_N=\min_{j\in\Z}|h+jN|$.  A geometric-progression sum, with ratio
$e_N(-hT_1)\ne1$, gives
\begin{equation}
\label{eq:hyperbola-geometric}
 |A(h)|,\ |B(h)|\ll\frac{N}{|hT_1|_N}\qquad(h\not\equiv0\pmod N).
\end{equation}
The right side is independent of the number of complete periods and
of the locations of the intervals.

Multiplication by $T_1$ permutes the nonzero residues modulo $N$ and
preserves their greatest common divisors with $N$.  Thus
\begin{equation}
\label{eq:hyperbola-gcd-sum}
\begin{split}
 \sum_{h\ne0\ (N)}\frac{(h,N)^{1/2}}{|hT_1|_N}
 &=\sum_{r\ne0\ (N)}\frac{(r,N)^{1/2}}{|r|_N}\\
 &\le\sum_{d\mid N}d^{1/2}
       \sum_{\substack{r\ne0\ (N)\\d\mid r}}\frac1{|r|_N}\\
 &\ll\log(2N)\sum_{d\mid N}d^{-1/2}
 \le\tau(N)\log(2N).
\end{split}
\end{equation}
Indeed the inner sum is at most
$2d^{-1}\sum_{1\le j\le N/(2d)}j^{-1}$, with a harmless possible
overcount of the midpoint.  Without the gcd weight the corresponding
sum is $O(\log(2N))$.  In conjunction with
\eqref{eq:hyperbola-geometric}, this proves
\begin{equation}
\label{eq:hyperbola-fourier-sums}
\begin{gathered}
 \sum_{h\ne0\ (N)}|A(h)|\ll N\log(2N),\qquad
 \sum_{h\ne0\ (N)}|A(h)|(h,N)^{1/2}
       \ll N\tau(N)\log(2N),
\end{gathered}
\end{equation}
and the identical estimates with $B$ in place of $A$.

\paragraph{Completion with the unit restriction.}
Define
\[
 S(h,k)=
 \sum_{\substack{u\bmod N\\(u,N)=1\\u\equiv a_0\ (T_0)}}
       e_N(hu+kc'\bar u).
\]
Fourier inversion, with the multiplicities in
\eqref{eq:hyperbola-fourier-weights}, gives the exact identity
\begin{equation}
\label{eq:hyperbola-completion-identity}
 C_l=\frac1{N^2}\sum_{h,k\bmod N}A(h)B(k)S(h,k).
\end{equation}
For example, summing $A(h)e_N(hu)/N$ over $h$ counts all allowed
integers $v\equiv u\pmod N$, which verifies the normalization directly.

Orthogonality modulo $T_0$ also gives the exact formula
\begin{equation}
\label{eq:hyperbola-restricted-sum}
 S(h,k)=\frac1{T_0}\sum_{j\bmod T_0}
       e_{T_0}(-ja_0)\Kl(h+jH,kc';N).
\end{equation}
The average over $j$ does not cost a factor $T_0$ in the following
absolute bounds.  When $k\ne0\pmod N$, the fact that $(c',N)=1$
and \eqref{eq:hyperbola-weil} give
\begin{equation}
\label{eq:hyperbola-second-nonzero}
 |S(h,k)|\ll\tau(N)\sqrt N\,(k,N)^{1/2},
 \qquad k\ne0\pmod N,
\end{equation}
uniformly also when $h=0$.  For $h\ne0$ and $k=0$, observe that
\begin{equation}
\label{eq:hyperbola-shifted-gcd}
 (h+jH,N)\le T_0(h,N).
\end{equation}
In fact any common divisor of $h+jH$ and $N$ divides $T_0h$, because
$T_0H=N$, and $(T_0h,N)\le T_0(h,N)$.  Consequently
\begin{equation}
\label{eq:hyperbola-first-nonzero}
 |S(h,0)|\ll\tau(N)\sqrt{NT_0}\,(h,N)^{1/2},
 \qquad h\ne0\pmod N.
\end{equation}
If a shift $h+jH$ vanishes modulo $N$, it is still covered by
\eqref{eq:hyperbola-shifted-gcd}; no shifted zero frequency is omitted.
No squarefreeness of $H$ or $N$ has been used.

At the two original zero frequencies,
\begin{equation}
\label{eq:hyperbola-complete-zero}
 S(0,0)=\frac{\varphi(N)}{\varphi(T_0)},\qquad
 \varphi(N)=T_0\varphi(H).
\end{equation}
For the first equality, reduction from the units modulo $N$ onto
the units modulo $T_0$ is surjective with equal fibers.  Explicitly,
if $t^a$ exactly divides $H$ and $t\mid T_0$, then the local modulus
in $N$ is $t^{a+1}$, and every unit modulo $t$ has $t^a$ unit lifts.
At a prime of $H$ outside $T_0$ there is no prescribed local class.
The Chinese remainder theorem proves the first equality.  The same
prime-power calculation gives the second.

\paragraph{The four frequency cases.}
The contribution of $h=k=0$ to
\eqref{eq:hyperbola-completion-identity} is
\begin{equation}
\label{eq:hyperbola-main-term}
 \frac{A(0)B(0)\varphi(N)}{N^2\varphi(T_0)}
 =\frac{UV\varphi(N)}{T_1^2N^2\varphi(T_0)}
   +O\left(\frac{U+V+1}{N}\right),
\end{equation}
where \eqref{eq:hyperbola-zero-weights} handles every endpoint
convention.  This is the stated main term since $UV=XY_2/l$.

For $h\ne0$, $k\ne0$, use
\eqref{eq:hyperbola-second-nonzero} and
\eqref{eq:hyperbola-fourier-sums}.  The total absolute contribution is
\[
 \ll\frac{\tau(N)\sqrt N}{N^2}
       \bigl(N\log(2N)\bigr)
       \bigl(N\tau(N)\log(2N)\bigr)
 =\tau(N)^2\sqrt N\log^2(2N).
\]
For $h=0$, $k\ne0$, the same argument with the zero weight retained
gives
\[
 \ll A(0)\tau(N)^2N^{-1/2}\log(2N).
\]
For $h\ne0$, $k=0$, instead use
\eqref{eq:hyperbola-first-nonzero}; this contribution is
\[
 \ll B(0)\tau(N)^2T_0^{1/2}N^{-1/2}\log(2N).
\]
Finally, $A(0)\ll U/T_1+1$ and $B(0)\ll V/T_1+1$.
These three estimates and the rounding term in
\eqref{eq:hyperbola-main-term} are all bounded by the error displayed
in \eqref{eq:hyperbola-explicit-count}.  This proves the general
assertion, with an absolute implied constant.

\paragraph{The size range needed in the variance argument.}
Under \eqref{eq:hyperbola-application-range},
\[
 R^{4/5}<N\le R^{1+o(1)},\qquad Y_2\le3N.
\]
The elementary divisor bound $\tau(N)=N^{o(1)}$ and $T_0=R^{o(1)}$
turn the explicit error into
\[
 R^{o(1)}\left(\sqrt N+\frac{X/l}{\sqrt N}
                           +\frac{Y_2}{\sqrt N}\right)
 \ll R^{1/2+o(1)}+R^{3/5+o(1)}
 =R^{3/5+o(1)}.
\]
The estimate uses only $l\ge1$, and is therefore uniform throughout
the asserted sieve range.  It also holds for shorter blocks and for
arbitrarily translated intervals.  This proves
\eqref{eq:hyperbola-application-count} and completes the proof.
\end{proof}

\begin{remark}
\label{rem:hyperbola-small-patterns}
In the application, $R\ge z^{\alpha_{\min}}$ with fixed
$\alpha_{\min}>0$,
$\mathcal Z=\exp(L/\log L)$, and
\[
 T_0=\prod_{\substack{t\le w\\t\mid H}}t,\qquad
 T_1=\prod_{\substack{t\le w\\t\nmid H}}t.
\]
Then $\log(T_0T_1)=O(w)=o(\log R)$ and
$\log\mathcal Z=o(\log R)$, uniformly in the isolated bin.
Thus the uniform family conditions in
\eqref{eq:hyperbola-application-range} follow from the parameters of
the main argument.  A fixed bin-width parameter may affect the point
at which these asymptotics apply, but it does not enter the absolute
constant in \eqref{eq:hyperbola-explicit-count}.
\end{remark}
